% ===================== 第 5 章 =====================
\section{三层记忆与 MemoBrain 适配}

\subsection{记忆是公共认知层，不是私有缓存}

这是多智能体推广带来的第一处变化。MemoBrain 的记忆是\textbf{私有}的：一个智能体写、同一个智能体读，压缩错了自己承担后果。议会里的记忆是\textbf{公共}的：一个席位写下的东西会被其他六个席位和研究者读到，所以压缩不再只是内部优化，而是一次\textbf{必须留痕的编辑}——这正是 Flush 要变成 Quarantine、Fold 要变成 Dialectical Fold 的原因。

Poliscope 把「记得什么」分成三层，各层回答不同的问题，也遵守不同的写入规则。

\begin{table}[ht]
\centering\small
\caption{三层记忆的分工}
\label{tab:memory}
\begin{tabularx}{\linewidth}{@{}L{2.8cm} L{3.2cm} X@{}}
\toprule
\rowcolor{accent}
\tblhead 层 & \tblhead 回答的问题 & \tblhead 内容与规则 \\
\midrule
私人记忆 Private Memory & 这位科学家怎么走到这一步 & 个人探索轨迹、工具记录、未验证猜想；只服务本人，他人不可读 \\
集体执行记忆 Collective MemoBrain & 共同体研究到哪一步、下一步查什么 & 结构化认知事件：谁提出什么、谁挑战谁、哪些争议悬着、哪些缺口没人碰；参与调度下一轮 \\
证据图 Evidence Graph & 我们目前对问题知道什么 & 只存经过审计的正式知识，唯一投影器可写（第 6 章） \\
\bottomrule
\end{tabularx}
\end{table}

三层之间是闭环而非堆叠：过程产生证据，证据暴露盲点，盲点作为新任务回灌私人记忆，驱动下一轮调查。集体记忆在这个闭环里承担执行控制职责——是证据状态决定下一步派谁查什么，而不是主持人按固定台本点名。

\subsection{MemoBrain 适配：小接口、不改上游}

MemoBrain 是外部方法基座（\code{github.com/qhjqhj00/MemoBrain}），集成前已核验许可证。Poliscope 通过 \code{MemoBrainAdapter} 对接，不修改上游源码，首版适配面刻意压到五个方法：

\begin{lstlisting}[language=ps]
init_private_memory(agent_id, task)
memorize_episode(agent_id, episode)
recall_private(agent_id, token_budget)
save_snapshot(agent_id)
load_snapshot(agent_id)
\end{lstlisting}

\begin{honestbox}
\boxtag{\color{warnamber}实现现状（不夸大）}
当前 \code{create\_memory\_adapter()} 返回的是 \code{InMemoryMemoryAdapter} 内存替身，过程图的 Flush/Fold/Recall 为简化实现，尚未真正导入上游 MemoBrain 代码。接口与边界已冻结，替换为真实实现不需要改动议会与证据层。这是首版明确保留项，详见第 12 章。
\end{honestbox}

\subsection{席位私有记忆的隔离规则}

隔离只靠一条规则做到位：agent id 由任务 id 与席位拼接，代码中不存在读取「另一个席位 id」的路径。于是同一任务的两席互相看不到对方的召回内容，同一席位也不会把 A 任务的记忆带进 B 任务。\code{CouncilMemory} 负责播种任务、记录 episode、按席召回与快照恢复。

\textbf{召回预算刻意设得很小：每席 800 字符}（\code{DEFAULT\_RECALL\_BUDGET}）。召回内容与证据投影竞争同一份上下文，它的目标是保住科研骨架而非重放完整转录——当前任务、已确认发现、活跃盲点、未决质询、少数异议、下一步证据需求，这六项在 800 字符内足以保全，而后面每一轮的提示词不会被不断增长的转录撑爆。

\subsection{过程记忆的 Fold：保骨架，不保字数}

\code{packages/memory/fold.py} 的 \code{fold\_process} 只接受过程图快照；任何试图折叠证据图节点的调用直接抛 \code{GraphBoundaryViolation}。压缩按任务、工具调用、决策三类节点进行，压缩前必须通过六项骨架保留校验：

\begin{center}
\begin{tcolorbox}[colback=panel, colframe=linegray, boxrule=0.6pt, arc=2pt, fontupper=\sffamily\small]
当前任务\enspace·\enspace 已确认 Finding\enspace·\enspace 活跃 Blindspot\enspace·\enspace 未解决质询\enspace·\enspace 少数异议\enspace·\enspace 下一步证据需求
\end{tcolorbox}
\end{center}

六项有一项丢失，折叠结果即标记为 rejected、原快照保留。

这里有一条容易忽略的分界：普通 Fold 只能作用于过程图。证据图只允许保留异议的辩证折叠（Dialectical Fold，第 6.7 节），两条折叠路径不可混用——它们保护的东西不同，一条守着上下文预算，另一条守着异议本身。

\subsection{Perspective Recall：同一份证据，不同的排序}

不同席位面对同一份证据快照时，按各自的 \code{evidence\_sort\_weights} 得到不同排序：因果专家先看到混杂与反向因果攻击，测量专家先看到构念冲突与测量不变性质疑。形式化地：
\[
\text{Context}_i = \text{ProcessRecall}_i + \text{EvidenceProjection}_i
\]
这同时避开了多智能体两个常见失败：每个 Agent 重复阅读全部材料，以及所有 Agent 最终收敛出相同观点。

去重的是检索成本，不是认知视角。

\subsection{认识论路由：让证据状态调度调查}

证据图发现缺口后，路由依次执行：生成盲点候选并打分 $\rightarrow$ 分派给全部七席（每席一个固定角色角度） $\rightarrow$ 分配角色化上下文 $\rightarrow$ 执行检索验证 $\rightarrow$ 结果写回集体记忆 $\rightarrow$ 重算认知前沿。注意这不是「选一个最强科学家单干」，也不是「一人认领一条」：七席全程在场，同一条盲点由七个角色角度同时调查，分工表在盲点悬赏轮\textbf{确定性}生成——同样的输入重跑一遍，得到的是同一张表。
